% Section 5: Applications
\subsection{Applications}
\label{sec:forecasting_rl_applications}

We outline four application domains where the prediction-and-decision interaction
is the active research question. The unifying claim across them is that joint or
decision-aware estimation outperforms two-step pipelines whenever the forecast
distribution is misspecified or the decision is asymmetric.

The contextual newsvendor is the canonical example. \citet{BanRudin2019} show that
on a high-dimensional feature space, an empirical risk minimisation procedure with
$\ell_1$ regularisation, applied directly to the newsvendor loss, dominates a
two-step pipeline that first estimates a demand distribution and then solves the
critical-fractile problem. On UK nurse staffing data, their single-step approach
reduces cost by between $23\%$ and $24\%$ relative to the sample average
approximation baseline, corresponding to roughly $\pounds 44{,}000$ to
$\pounds 47{,}000$ per ward per year, and their finite-sample bounds extend the
\citet{Levi2007} guarantees to the high-dimensional contextual case. The
\citet{Oroojlooyjadid2020} neural extension parametrises the order quantile
function directly with a deep network and trains it on the asymmetric newsvendor
loss. \citet{HanZhengMisra2024} introduce a conditional deep generative model that
respects the decision-dependent endogeneity in joint demand learning and pricing,
proving consistency for the joint problem. \citet{Amazon2025} embed a probabilistic
drain forecaster as a differentiable simulator inside an RL inventory ordering
loop and validate the system on counterfactual out-of-distribution inventory
states that arise from off-policy trajectories.

Electricity dispatch is the cleanest empirical decision-focused win.
\citet{DontiAmosKolter2017} on PJM grid data already established that task-based
training reduces realised cost over likelihood-based training when the same
optimisation routine is used downstream. \citet{CaoXu2025} extend this to microgrid
operation by embedding a multilayer probabilistic encoder-decoder forecaster
inside a two-stage robust optimisation, training under a smart predict-then-optimize
surrogate. On IEEE 33-bus and 69-bus benchmarks, they report total cost reductions
of up to $18\%$ relative to two-stage pipelines and improved forecast accuracy
along the operational frontier. Their phrasing is the cleanest statement of the
chapter's thesis, namely that the proposed method aligns forecasting accuracy
with operational objectives by directly minimising decision regret via a surrogate
SPO loss function.

Macroeconomic nowcasting and monetary policy reaction sit at the intersection of
our two halves. \citet{Fernandes2024} demonstrates the forecasting half. The
decision-side example is \citet{HinterlangTaenzer2024}. They estimate a structural
vector autoregression on US data from 1987Q3 to 2007Q2 and use it as the
transition model for a reinforcement learning agent that learns a nonlinear
monetary policy reaction function under a quadratic central bank loss. The
learned policies outperform the realised federal funds rate path and standard
Taylor rules on the calibrated loss, and the framework accommodates the zero
lower bound and nonlinear preferences without requiring closed-form solution
of the optimal policy problem.

Order execution and market microstructure provide a fourth application set.
\citet{MicheliMonod2024} train a Deep Deterministic Policy Gradient agent for
optimal execution under transient price impact, using an auxiliary $Q$-function
that projects the non-Markovian state into a low-dimensional representation. The
agent converges to the optimal trading rate across exponential, power-law, and
linear resilience kernels, and adapts online when the kernel decay parameter
drifts linearly in time. \citet{Briola2023} train an Apex-Dueling-DQN agent on
NASDAQ Apple limit-order-book data with a noisy directional forward-mid-quote
signal and report return uplifts of between $14.8$ and $32.2$ percentage points
relative to a heuristic benchmark with access to the same signal, with the gains
robust to the 2012 to 2022 regime shift at low and moderate signal noise. The
forecast is the noisy signal, and the agent's contribution is to translate the
signal into a profitable trading rule that the heuristic does not extract.
